\chapter{System Overview}
\label{sec:overview}

This chapter gives a high-level view of the system before the implementation details appear. The emphasis is on workflow, component roles, and evidence flow rather than source-tree organization.

\section{End-to-End Workflow}

A run begins with a user query or evaluation task, resolves a workflow configuration, provisions tool access, and then executes a staged analysis pipeline. The important point is not the exact call chain but the way configuration, tooling, and staged reasoning are assembled into one repeatable workflow. Appendix~\ref{app:reproduction-commands} provides the concrete launch details for readers who need the artifact-level execution path.

In this thesis, a \emph{pipeline} is the ordered sequence of stages used for one run, and a \emph{stage} is a bounded phase with its own prompt contract, tool access pattern, and expected output. The repository also uses the term \emph{stage manager} for the coordinator prompt and runtime logic attached to a single stage. Likewise, a worker \emph{archetype} is a reusable specialist role definition such as \texttt{ghidra\_analyst}, while a \emph{subagent} is a concrete worker instance or delegated model invocation created from one of those archetypes during a run.

In the maintained \texttt{auto\_triage} baseline used in the final experiments, execution begins with a light preparatory pass, called \emph{preflight} in the system implementation, and deterministic presweeps that gather reusable context before any broad free-form reasoning occurs. A planner then converts the request and shared context into bounded work items, specialist workers execute those items with tool access appropriate to their role, a dedicated proposal-writing stage normalizes conservative change suggestions into queued proposals, and a reporter synthesizes the result into the final analysis artifact. Some experimental pipelines insert an explicit validator stage before final reporting, but validator review is not universal across all runs. After the final report is produced, the harness captures artifacts, invokes the configured LLM judge, and aggregates results into the experiment record.

\section{Planner, Worker, Proposal-Writer, Validator, and Reporter Roles}

The workflow supports multiple stage kinds, but its core pattern is straightforward: prepare context, decompose the task, gather evidence, optionally review coverage, and then synthesize the final answer.

Preparatory stages establish basic sample metadata and reusable static-analysis outputs. Planners decide how to decompose the task and which evidence should be gathered. Workers perform focused analysis passes, often with specialized tool access or role prompts that constrain their scope. Proposal-writing stages, when present, convert worker-surfaced rename, type, comment, or patch suggestions into structured queued changes rather than leaving that burden to the reporter. Validators, when present, act as an internal quality gate by reviewing coverage, evidence quality, and missing artifacts before final reporting. Reporters convert intermediate evidence into the final narrative artifact used by the harness and LLM judge.

\section{Configuration Surfaces}

The system exposes several configuration families that are central to the evaluation:
\begin{itemize}
\item \textbf{Response-scope variants} control how concise or expansive the final report should be.
\item \textbf{Worker subagent profiles} determine whether workers are instantiated as generalists or as specialist archetype-driven subagents.
\item \textbf{Worker prompt shape and role-prompt mode} tune the instruction surface given to each worker.
\item \textbf{Tool profiles} control which MCP tools are available to the pipeline.
\item \textbf{Architecture presets} define the worker mix.
\item \textbf{Pipeline presets} select the stage sequence.
\item \textbf{Validator topology and review level} determine how strict the internal quality gate is.
\end{itemize}

These are not merely engineering options; they are the intervention points used by the evaluation harness to compare baseline and variant behavior. Some of them primarily change how work is organized, such as worker subagent profiles, architecture presets, and pipeline presets. Others change how evidence is requested or summarized, such as response-scope variants and worker-prompt surfaces. Tool profiles matter most when asking how much performance depends on breadth of tool access rather than on orchestration alone. The full mapping of variant labels to configuration changes is provided in Table~\ref{tab:evaluation-variants}.

This distinction becomes important later in the thesis because these configuration surfaces are the main objects of RQ2. The experimental question is not whether the system has options, but whether changing those options materially affects analysis quality, task success, or evidence quality on the benchmark.

\section{Shared State and Evidence Flow}

The system maintains shared run state so that later stages can inspect earlier findings, tool outputs, and failure metadata. This matters because reverse engineering benefits from explicit provenance: shared state determines what later stages can trust, what evidence survives handoff between agents, and what ultimately appears in final reports and judge inputs.

Shared state captures four categories of information: intermediate tool outputs and structured findings from earlier stages; stage status logs and work-item tracking; the final analysis report and supporting artifacts; and run metadata used by the evaluation harness, including bundle provenance and comparison context.

That evidence flow is also central to RQ3. If artifacts are not preserved clearly enough to survive handoff from workers to validators and reporters, then the system may still produce a plausible report while failing the more important question of whether its claims remain visibly grounded.

\section{Why This Architecture Fits Reverse Engineering}

The same shared-state model also supports the two user-facing surfaces discussed later in the thesis. The analyst-facing interface exposes pipeline progress, tool logs, and intermediate stage outcomes while a run is in progress. The evaluation-facing browser exposes archived runs, repaired comparisons, and per-executable result views after the fact. Those interfaces matter scientifically because they make provenance inspectable rather than leaving the workflow as a single opaque report.

The value proposition of the architecture is straightforward: binary analysis benefits from staged specialization because the evidence sources are heterogeneous, and high-confidence answers require both discovery and disciplined synthesis. A monolithic prompt can mention these needs, but it cannot enforce them structurally. The architecture used here makes those roles explicit so they can be inspected, varied, and measured.
